\section{Theoretical Analysis}
\label{sec:theory}

\subsection{Feasible-space restriction}

\begin{proposition}
For every contiguous platoon partition \(\Pi\), the platoon-constrained optimality gap is nonnegative:
\begin{equation}
G(\Pi)\ge 0.
\end{equation}
\end{proposition}

\begin{proof}
Every sequence that satisfies the platoon-continuity constraints also satisfies the original FIFO constraints. Hence \(\pseqset\subseteq\seqset\). Minimization over a subset cannot produce an objective value smaller than minimization over the original set.
\end{proof}

\subsection{Global sequence repair}

The analysis does not decompose the system into independent one-platoon examples. Instead, it starts from a complete vehicle sequence and repairs every violated platoon adjacency while retaining all interactions with other approaches.

Consider two consecutive vehicles \(a\) and \(a'\) from the same approach and platoon. Suppose a current FIFO-feasible sequence has the form
\begin{equation}
s=[P,a,B,a',R],
\end{equation}
where \(B=(b_1,\ldots,b_m)\) is a nonempty sequence containing vehicles from any number of other approaches. The repaired sequence is
\begin{equation}
s'=[P,a,a',B,R].
\end{equation}
Let \(d=r_{a'}-r_a\), \(q=\max\{d,\hF\}\), and let \(M=m+|R|+1\) be the number of vehicles whose passing times may change.

\begin{lemma}[Candidate local perturbation bound]
\label{lem:local-repair}
Under the assumptions in Section~\ref{sec:formulation}, the change in total delay caused by the repair satisfies
\begin{equation}
J(s')-J(s)
\le
M\pospart{d-\hF}-2(\hS-\hF),
\label{eq:candidate-local-bound}
\end{equation}
where \(J(s)=ND(s)\).
\end{lemma}

\statuscandidate

The complete draft proof, including every comparison used in the recursion, is contained in \texttt{notes/proof\_development.tex}. Equation~\eqref{eq:candidate-local-bound} must be exhaustively tested before submission.

\subsection{Candidate global bounds}

Starting from an original optimal sequence, repeatedly apply the repair operation until every platoon is contiguous. A telescoping argument gives the following candidate sequence-specific result:
\begin{equation}
G(\Pi)
\le
\frac{1}{N}
\sum_{e\in\mathcal{Q}}
\pospart{M_e\pospart{d_e-\hF}-2(\hS-\hF)},
\label{eq:candidate-sequence-bound}
\end{equation}
where \(\mathcal{Q}\) is the set of internal links that require repair and \(M_e\) is evaluated at the corresponding repair step.

Using \(M_e\le N\) produces the partition-specific candidate bound
\begin{equation}
B(\Pi)
=
\sum_{e\in\mathcal{E}_{\Pi}}
\pospart{
\pospart{d_e-\hF}
-\frac{2(\hS-\hF)}{N}
},
\qquad
G(\Pi)\le B(\Pi).
\label{eq:candidate-partition-bound}
\end{equation}

If every platoon internal gap satisfies \(d_e\le\delta\), then
\begin{equation}
G(\Pi)
\le
(N-K)
\pospart{
\delta-\hF-\frac{2(\hS-\hF)}{N}
}.
\label{eq:candidate-parameter-bound}
\end{equation}

\statuscandidate

Before these bounds are promoted to final theorems, the study will: (i) enumerate all FIFO sequences and contiguous partitions for small systems; (ii) search for counterexamples; (iii) measure the ratio \(G(\Pi)/B(\Pi)\); and (iv) identify the inequalities responsible for conservativeness.
